% Use case (short): symmetry-resolved entanglement in the XXZ chain. Full account: man/btc_sre (Part B).
% Numbers: scripts/symmetry_resolved.py -> generated/symmetry_resolved.tex (macros \SR...).
\subsection{Use case: symmetry-resolved entanglement in the XXZ chain}\label{sec:sre}

For a pure state invariant under a $U(1)$ symmetry, the reduced state $\rho_A$ commutes with the charge $Q_A$ of the
subsystem and splits into sector blocks, $\rho_A=\bigoplus_qp(q)\rho_A^{(q)}$, and $S(\rho_A)=\sum_qp(q)S(q)+H(p)$
\cite{laflorencie2014,goldstein2018,xavier2018}. In critical chains the resolved entropies $S(q)$ are equipartitioned to
leading order~\cite{xavier2018}; the corrections and their asymptotics for free fermions are due to Bonsignori, Ruggiero
and Calabrese~\cite{bonsignori2019}, and the variance of $Q_A$ grows as $(K/\pi^2)\log\ell$ with the Luttinger parameter
$K$~\cite{song2012}. In the library the reduced state of a charge-conserving state \emph{is} an operator of
$\bigoplus_qM_{d_q}$, so the resolved quantities are read off its blocks without projecting.

For the XXZ chain at $L=\IfFileExists{generated/symmetry_resolved.tex}{\SRL}{24}$, the ground states for eight anisotropies
come from one batched Lanczos run in the sector $S^z=0$, and the resolved entropies, the charge statistics and the widths
$B_n$ of the charge distributions of $\rho_A^n$ follow for all subsystem sizes (Figure~\ref{fig:sre}). The charge
statistics give the Luttinger parameter to within
\IfFileExists{generated/symmetry_resolved.tex}{$\SRKFcsDevPct\,\%$ for $\SRKFcsMidRange\le\Delta\le\SRKFcsMidRangeHi$}{about
one per cent} (about $4\%$ at $\Delta=-0.75$, and with the known logarithmic corrections at the isotropic point); the
$1/n$ law of the widths $B_n$ is tested and the amplitude $G_1(\Delta)$ of the breaking of equipartition measured; for free
fermions the asymptotic value $1/(2\pi^2)$ of~\cite{bonsignori2019} is recovered and its slow approach shown. Details,
validation and tables are in the companion paper~\cite{nn_btc}.

\begin{figure}[h]
\centering\fig{symmetry_resolved.pdf}{\textwidth}
\caption{Periodic chain, $L=\SRL$. Left: $nB_n$ against $\log\ell_c$, $n=1$ (filled circles), $2$ (squares), $3$
(triangles); lines: slope $K/\pi^2$ (Bethe ansatz). Middle: $K$ from the charge statistics and from the widths, against
the Bethe-ansatz value. Right: $S(q)-S(\ell/2)$ at $\ell=\SRHalf$ for $\Delta\in\{\SRDeltas\}$ (dark to light); dotted: the
Gaussian form, exact at $q-\ell/2=\pm1$.}
\label{fig:sre}
\end{figure}
